\documentclass[10pt]{article}

\usepackage[utf8]{inputenc}

\usepackage[T1]{fontenc}

\usepackage{amsmath}

\usepackage{amsfonts}

\usepackage{amssymb}

\usepackage[version=4]{mhchem}

\usepackage{stmaryrd}

\usepackage{bbold}



\begin{document}

ной Г. г. в теории линейных дифференциальных уравнений аналогична роли обычной группы Галуа в теории алгебраич. уравнений. Пусть $k-$ дифференциальное подполе поля мероморфных функций от $z$ с дифференцированием $\frac{d}{d z}$, $\bar{k} \supset \mathbb{C}$, а $K$ порождается над $k$ элементами $y_{i j}$ фундаментальной матрицы решений уравнения $y^{\prime}=A y, A \in \operatorname{Mat}(n, k)$. Тогда: 1) возможность выразить $y_{i j}$ через элементы $k$ с помощью операций,$+ \times$, взятия первообразной, взятия экспоненты первообразной и решения алгебраич. уравнений эквивалентна разрешимости связной компоненты единицы группы $G$; 2) $G$ - алгебраич. подгруппа в $G L(n, \mathbb{C})$, то есть $G$ выделяется в $G L(n, \mathbb{C})$ уравнениями, полиномиальными относительно матричных элементов; 3) если $k-$ поле рациональных функций, а особые точки уравнения $y^{\prime}=A y$ регулярны, то $G-$ наименышая алгебраич. подгруппа, содержащая монодромии группу уравнения.



Лит.: [1] К а план с к и й И., Введение в дифференциальную алгебру, пер. с англ., М., 1959; [2] K ol c h i n E. R., Differential algebra and algebraic groups, N.Y., 1973; [3] Pic ard E., Traité d'analyse, t. 3, $\begin{array}{ll}\text { ch. 17, Р., } 1896 . & \text { В. Г. Дринфельд. } \\ & \end{array}$ ГАЛЬТОНА-ВАТСОНА ПРОЦЕСС - см. Ветвящийся проuecc.



ГАМИЛЬТОНА КАНОНИЧЕСКОЕ УРАВНЕНИЕ - УРавнение



\[

\dot{q}=H_{p}^{\prime}, \quad \dot{p}=-H_{q}^{\prime},

\]



описывающее движение динамической системы с гамильтонианом $H(p, q, t)$ в координатах Дарбу. $\quad$ М. Э. Казарян. ГАМИЛЬТОНА КВАДРАТИЧНАЯ ФУНКЦИЯ - сМ. Квадратичный гамильтониан.



ГАМИЛЬТОНА OПEPATOP - см. Гамильтониан.



ГАМИЛЬТОНА УРАВНЕНИЯ, каноническ ие урав нения механики, - дифференциальные уравнения движения голономной механической системы в канонических переменных, которыми являются $s$ обобщенных координат $q_{i}$ и $s$ обобщенных импульсов $p_{i}$, где $s-$ число степеней свободы системы. Выведены У. Гамильтоном (W. Hamilton) в 1834. Для составления Г. у. надо в качестве характеристич. функции системы знать Гамильтона функцию $H\left(q_{i}, p_{i}, t\right)$, где $t$ - время. Тогда, если все действующие на систему силы потенциальны, Г.у. имеют вид



\[

\frac{d q_{i}}{d t}=\frac{\partial H}{\partial p_{i}}, \quad \frac{d p_{i}}{d t}=-\frac{\partial H}{\partial q_{i}} .

\]



Если наряду с потенциальными на систему действуют непотенциальные силы $\boldsymbol{F}$, то к правым частям 2-й группы уравнений $\left({ }^{*}\right)$ надо прибавить соответствующие обобщенные силы $Q_{i}$. Уравнения $\left({ }^{*}\right)$ представляют собой систему $2 s$ обыкновенных дифференциальных уравнений 1-го порядка, интегрируя к-рые можно найти все $q_{i}$ и $p_{i}$ как функции времени $t$ и $2 s$ постоянных интегрирования, определяемых по начальным данным. Решение системы уравнений $\left({ }^{*}\right)$ можно также свести к отысканию полного интеграла соответствующего ей уравнения с частными производными (см. Гамильтона-Якоби уравнение ).



Если одна из координат $q_{i}$, напр. $q_{1}$, является циклич. координатой, то есть явно не входит в выражение функции $H$, то $\partial H / \partial q_{1}=0$ и одно из уравнений $\left({ }^{*}\right)$ дает сразу интеграл $p_{1}=\alpha_{1}$, где $\alpha_{1}-$ постоянная. Особый интерес представляет случай, когда все координаты циклические, а функция $H=H\left(p_{i}\right)$ явно не зависит от времени (силовое поле и наложенные связи стационарны). Тогда все $p_{i}=\alpha_{i}$, то есть постоянны; следовательно, функции $H\left(p_{i}\right)$ и $\partial H / \partial p_{i}$ тоже постоянны, и 1-я группа уравнений $\left(^{*}\right)$ дает $d q_{i} / d t=\beta_{i}$, откуда $q_{i}=\beta_{i} t+C_{i}$, где $\beta_{i}, C_{i}$ - новые постоянные. Уравнения в этом случае интегрируются элементарно и все координаты являются линейными функщиями времени. Отсюда следует, что задачу интегрирования Г.у. можно свести к



110 ГАМИЛЬТОНА задаче отыскания для системы циклич. координат. Это, в принципе, возможно, так как Г.у. обладают тем важным свойством, что они допускают переход с помощью так наз. канонич. преобразований от переменных $q_{i}, p_{i}$ к новым переменным $Q_{i}\left(q_{i}, p_{i}, t\right), P_{i}\left(q_{i}, p_{i}, t\right)$, к-рые также являются каноническими и удовлетворяют уравнениям $\left({ }^{*}\right)$ с соответствующей функцией $H\left(Q_{i}, P_{i}, t\right)$.



Равноправность в Г.у. координат и импульсов как независимых переменных, а также инвариантность этих уравнений по отношению к канонич. преобразованиям открывают большие возможности для обобщений. Поэтому Г. у. имеют важные приложения не только в механике, но и во многих других областях физики, напр. в статистич. физике, кванто$\begin{array}{ll}\text { вой механике, электродинамике и др. } & \text { C. М. Тарг. }\end{array}$ ГАМИЛЬТОНА ФУНКция - характеристическая функция механической системы, выраженная через канонические переменные: обобщенные координаты $q_{i}$ и обобщенные импульсы $p_{i}$. Для системы со связями, явно не зависящими от времени $t$, движущейся в стационарном потенциальном силовом поле, Г.ф.



\[

H\left(q_{i}, p_{i}\right)=T+U,

\]



где $U$ - потенциальная, а $T-$ кинетич. энергия системы, в выражении к-рой произведена замена всех обобщенных скоростей $\dot{q}_{i}$ на $p_{i}$ с помощью равенств



\[

p_{i}=\partial T / \partial \dot{q}_{i} .

\]



Таким образом, Г. ф. равна в этом случае полной механич. энергии системы, выраженной через $q_{i}$ и $p_{i}$. В общем случае Г. ф. $H\left(q_{i}, p_{i}, t\right)$ может быть определена через Лагранжа функцию $L\left(q_{i}, \dot{q}_{i}, t\right)$ равенством



\[

H\left(q_{i}, p_{i}, t\right)=\left[\sum_{i} p_{i} \dot{q}_{i}-L\left(q_{i}, \dot{q}_{i}, t\right)\right]_{\dot{q}_{i} \rightarrow p_{i}},

\]



в к-ром все $\dot{q}_{i}$ должны быть выражены также через $p_{i}$.



Г. ф., как и функция Лагранжа, полностью характеризует ту систему, для к-рой она определена, так как, зная $H\left(q_{i}, p_{i}, t\right)$, можно составить дифференциальные уравнения движения системы или в виде $2 s$ обыкновенных дифференциальных уравнений 1-го порядка, где $s-$ число степеней свободы, или в виде одного уравнения с частными производными тоже 1-го порядка (см. Гамильтона-Якоби уравнение). Г. ф. введена У. Гамильтоном (W. Hamilton).



Наряду с термином «Г. ф.» употребляют иногда термин «главная функция Гам ильтона», именуя так полный интеграл уравнения Гамильтона - Якоби, равный действию по Гамильтону. В квантовой механике используется квантовомеханич. оператор - гамильтониан, или оператор Гамильтона, соответствующий Г. ф. в классич. механике.



C. M. Tap2.



ГАМИЛЬТОНА ФУНКЦИЯ, га м и л ь то н и а н,- произвольная гладкая функция на симплектическом многообразии $(M, \omega)$. Г. ф. Н задает динамич. систему на $M$ с помошью векторного поля $\boldsymbol{v}_{H}$, определяемого по правилу:



\[

\omega\left(\cdot, v_{H}\right)=d H .

\]



Система дифференциальных уравнений, описывающая гамильтонову систему, в координатах Дарбу принимает вид Гамильтона канонического уравнения



\[

\dot{p}=-H_{q}^{\prime}, \quad \dot{q}=H_{p}^{\prime} .

\]



М. Э. Казарян.



ГАМИЛЬТОНА ФУнКЦИЯ; нормальная формапростейший вид, к которому можно привести Гамильтона функцию $H$ в окрестности: а) положения равновесия, б) замкнутой траектории с помощью канонической замены координат [в случае a) предполагается, что $H$ не зависит или периодически зависит от времени, а в случае б) - что $H$ не зависит от времени].



Те орема. Пусть в окрестности точки $O \in \mathbb{R}^{2 N}$ задана $C^{\infty}$-гладкая Г. ф.



\[

H=H(p, q)=O\left(|p|^{2}+|q|^{2}\right)

\]



причем все собственные числа $\pm i \omega_{1}, \ldots, \pm i \omega_{n}$ линеаризации в 0 соответствующего гамильтонова векторного поля чисто мнимы и нерезонансны, то есть $m_{1} \omega_{1}+\ldots+m_{n} \omega_{n} \neq 0$ для





\end{document}